\subsection{\resultrevision{Comparison Design}}
\begingroup
\color{resultrevision}
The main comparison places Full MPSLGP alongside the published PSLGP-I results in Tables~IV--VI of PSLGP~\cite{xu2026pslgp}. Both use IGD-based training and 30 evolutionary seeds. Full learns the task-specific conditioned heuristics jointly with 1,000 individuals per group; the published PSLGP-I experiments allocate 1,000 individuals to each independently trained task. The published entries are contextual comparisons. Their raw fronts and reference data are not available for recomputation, so the main comparison is not a common-reference statistical test of superiority.

No-transfer and fixed-transfer are completed ablations, not additional main baselines. Each has the same eight groups, seeds 0--29, one-tree crossover, task quotas, IGD training, surrogate procedure, and test configuration as Full. The verified no-transfer profile disables adaptive transfer and sets its transfer probability and budget to zero. Fixed-transfer changes only the adaptive-transfer flag relative to Full; its fixed mating probability is 0.8 for two-task groups and 0.4 for three-task groups. All non-ablation training fields match. The realized adaptive probabilities need not match these fixed rates.

We recompute HV, IGD, and GD for all three MPSLGP methods using one common set of objective bounds and an empirical reference front for each benchmark scenario, pooled across methods, groups, and seeds. Higher HV and lower distances are preferable. Invalid preferences are reported separately; control quality scores are conditional on valid preferences, and every training seed is retained. The statistical unit is a paired training seed for a fixed group/task target. Means and sample standard deviations are reported over 30 runs. Exact two-sided signed-rank tests omit zero differences and use midranks for absolute ties. Holm correction covers all 108 target/metric/control comparisons at family-wise level 0.05.

The completed ablations test transfer presence and adaptive versus fixed transfer under the recorded configurations. They do not individually identify the effects of preference regions, reward terms, or task inheritance. A budget- and operator-matched independent PSLGP rerun, a $K=1$ controller, and a frequency-matched fixed controller remain useful follow-up controls; their results are not claimed in this paper.
\endgroup